\documentclass[UTF8,fontset=none,11pt,a4paper]{ctexart}
\usepackage{ipara-handbook}
\handbooksetup{DS-I01 从 Workload 到数据结构选择}{408 · 数据结构 Integration I01}{Workload · Composition · Boundary}
\begin{document}
\handbooktitle{从 Workload 到数据结构选择}{把“会背结构名”升级为可审计、可反证的组合设计能力}{408 · 数据结构 Integration I01}

\begin{corebox}{完整组合问题：先写契约，再拼能力}
\[
\boxed{
\begin{aligned}
\text{Workload}
&\to \text{API Contract}
\to \text{Required Capabilities}
\to \text{Candidate Structures}\\
&\to \text{Composite State / Invariants}
\to \text{Cost Vector}
\to \text{Choice}
\to \text{Boundary Check}.
\end{aligned}}
\]
本册不重新教授任何单一结构，而是检查能否把 Atlas Foundation、Topic 和 Bridge 的接口组合成一项有理由、可反证、可迁移的设计。一个 API 往往要求多种能力；单一结构不能同时提供时，正确动作不是寻找“万能结构”，而是让多个 Owner 各自承担一种能力，并显式维护它们之间的同步不变量。
\end{corebox}

\begin{methodbox}{Position、Owns、Uses 与 Stop Boundary}
\begin{itemize}
  \item \kw{Position：}本册是数据结构 Integration，拥有“从需求到组合方案”的完整决策轨迹。
  \item \kw{Owns：}API/工作负载解析、能力分解、结构组合、跨结构同步不变量、规模升级路线、条件化选择与独立验证。
  \item \kw{Uses：}DS01 的数组/链表表示，DS05 的优先队列，DS09 的有序索引，DS10 的 Hash，DS12 的外部排序，以及 DS-B02 的 Workload--Index 比较接口。
  \item \kw{Stops：}Hash 冲突、堆上滤/下滤、链表改边、外部归并等局部机制回各 Topic；线程调度、JVM GC、MapReduce 执行引擎只作为工程 Extension，不在本册重定义。
\end{itemize}
若缺少操作频率、最坏保证、数据分布或资源约束，应输出条件化选择，而不是伪造唯一答案。
\end{methodbox}

\tableofcontents
\newpage

\section{先把自然语言需求改写成可检查的合同}

\subsection{API 名字不等于真正需求}
“设计缓存”“维护热搜”“做一个随机集合”仍然太粗。开始选结构前，应把每个操作写成五元组：
\[
\boxed{(\text{Input},\ \text{Output},\ \text{Mutation},\ \text{Frequency},\ \text{Required Bound})}.
\]
例如随机集合的合同不是“支持集合操作”，而是：
\begin{center}\small
\begin{tabularx}{\linewidth}{L{2.3cm}YYL{2.2cm}}
\toprule
操作 & 语义 & 状态变化 & 目标成本\\
\midrule
\code{insert(x)} & 若不存在则加入并返回真 & 增加一个可随机抽样元素 & 平均 $O(1)$\\
\code{remove(x)} & 若存在则删除并返回真 & 删除后其余元素仍可随机抽样 & 平均 $O(1)$\\
\code{getRandom()} & 对现存元素做等概率抽样 & 只读 & $O(1)$\\
\bottomrule
\end{tabularx}
\end{center}
“等概率”会生成数组的紧凑下标能力，“存在性与定位”会生成 Hash 的直接定位能力，“任意位置删除”则产生二者之间的交接压力。

\subsection{能力分解：每个成本承诺需要谁来兑现}
常见能力与默认候选如下：
\begin{center}\small
\begin{tabularx}{\linewidth}{L{3.0cm}L{3.3cm}YY}
\toprule
所需能力 & 常见 Owner & 能得到什么 & 必须付出什么\\
\midrule
按 key 直接定位 & Hash & 平均常数查找 & 冲突、装填、重散列、无序\\
按位置随机访问 & 紧凑数组 & $O(1)$ 下标访问 & 中间插删需移动\\
已知节点局部删除 & 双向链表 & $O(1)$ 改边 & 需要额外定位与指针空间\\
持续取得当前极值 & Heap / Priority Queue & 查看极值 $O(1)$、更新 $O(\log n)$ & 不提供任意排名和全序\\
范围、前驱、后继 & 有序索引 & 有序查询与扫描 & 更新需维护有序/平衡\\
前缀匹配 & Trie & 查询成本绑定 key 长度 & 节点稀疏与空间放大\\
超内存排序/聚合 & 分片、外部归并 & 把内存约束转成多阶段处理 & I/O、缓冲、临时文件与倾斜\\
\bottomrule
\end{tabularx}
\end{center}
这张表不是答案表。它只给候选；最终选择仍由操作比例、最坏保证、输入分布、内存、I/O 与是否允许近似共同决定。

\subsection{不要先追求最优：先写可证明的基线}
组合设计遵循三步：
\begin{enumerate}
  \item 写出语义正确的基线，例如全量排序、线性扫描或单一 Map；
  \item 指出基线重复了什么工作或违反了哪一项资源约束；
  \item 只增加能消除该瓶颈的状态，并写出新增状态的维护成本。
\end{enumerate}
“换成更高级的数据结构”不是推导。必须说明它预存了什么信息，为什么该信息足以减少重复工作，以及更新时怎样保持真实。

\section{组合结构的核心不是容器数量，而是同步不变量}

\subsection{一份事实可以有多个索引，但不能有多个真相}
组合结构通常保存一份逻辑对象和多个访问路径：
\[
\text{Logical Record}
\longleftrightarrow
\begin{cases}
\text{Key Index},\\
\text{Position / Order Index},\\
\text{Priority / Frequency Index}.
\end{cases}
\]
真正的难点是所有索引必须指向同一批有效对象。常见同步不变量包括：
\begin{enumerate}
  \item Map 中每个 key 恰好指向一个仍存活的节点或下标；
  \item 数组有效区间 $[0,n)$ 没有空洞，Map 中记录的下标与数组一致；
  \item 链表中每个真实节点恰好出现一次，头尾哨兵互相可达；
  \item 频率桶内的 key 与 key 表中的频率一致，空桶不参与最小频率判断；
  \item 两个堆的分区有序且大小差不超过一。
\end{enumerate}

\begin{warnbox}{最危险的局部更新}
组合结构中的 bug 往往不是某个容器操作错误，而是只更新了一个索引。例如数组尾元素搬到被删位置后若忘记更新 Map 中的下标，第一次删除看似成功，下一次定位就会访问错误对象。每个修改操作都应列出“哪些视图变了”，并在返回前逐一恢复不变量。
\end{warnbox}

\subsection{原子状态转移清单}
即使实现语言没有事务，也可以按事务思维审查：
\[
\begin{aligned}
\text{Validate Preconditions}
&\to \text{Identify All Affected Views}
\to \text{Mutate} \\
&\to \text{Repair Cross-Index Links}
\to \text{Check Size / Boundary}.
\end{aligned}
\]
若中途可能失败，还要决定回滚、延迟提交还是先构造新状态再交换。408 伪代码通常不要求异常安全，但顺序仍决定边界是否正确。

\section{生成性例一：Hash + Array 实现随机集合}

\subsection{为什么“交换到末尾再删除”是必要接口}
数组支持 $O(1)$ 随机下标，但任意中间删除会移动后缀。Map 能立即找到目标下标，却不能消除移动。解决方法是把“删除任意位置”翻译成“删除末尾”：
\[
\boxed{
\text{locate }i
\to A[i]\leftarrow A[n-1]
\to \text{repair moved element's index}
\to \text{pop tail}
\to \text{erase key}.}
\]
它牺牲元素顺序，换来紧凑数组与常数删除；若题目要求保持顺序，这个方案立即失效。

\begin{lstlisting}[language=C++,caption={随机集合的状态合同伪代码}]
state:
  values[0..n)            // dense, no holes
  pos[value] -> index     // exact inverse index

insert(x):
  if x in pos: return false
  pos[x] = n
  values.push_back(x)
  return true

remove(x):
  if x not in pos: return false
  i = pos[x]
  y = values[n - 1]
  values[i] = y
  pos[y] = i              // repair before/with shrinking
  values.pop_back()
  erase pos[x]
  return true
\end{lstlisting}
当 $x$ 本来就是末元素时，$y=x$；先写回 \code{pos[y]} 再删除 \code{pos[x]} 仍可得到正确结果。空集合上的随机抽样必须由 API 预条件或显式失败语义处理。

\section{生成性例二：Hash + 双向链表实现 LRU}

\subsection{能力分工}
LRU 要求 \code{get/put} 都为平均 $O(1)$：
\begin{itemize}
  \item Hash：由 key 直接找到节点；
  \item 双向链表：已知节点时 $O(1)$ 删除，并把它移到“最近使用”一端；
  \item 容量计数：超过 capacity 时从“最久未使用”一端淘汰。
\end{itemize}
单链表不能在仅持有当前节点时 $O(1)$ 找到前驱；节点必须保存 key，否则从链表淘汰尾节点后无法同步删除 Hash 项。

\subsection{状态与不变量}
设头哨兵之后是最近使用端，尾哨兵之前是最久未使用端：
\[
\boxed{
\begin{gathered}
\operatorname{keys}(Map)=\operatorname{keys}(List),\\
Map[k]\text{ 指向链表中唯一的 }k\text{ 节点},\\
\text{List order}=\text{nonincreasing recency},\\
|Map|\le capacity.
\end{gathered}}
\]
\code{get(k)} 命中不仅返回值，还会改变 recency，因此是逻辑写操作；\code{put} 更新旧 key 时也必须刷新其最近使用位置。容量为零是独立边界，任何插入都不能留下节点。

\subsection{为什么使用哨兵}
虚拟头尾把“空表、首节点、尾节点、中间节点”统一为同一种四指针改边，减少分支但不消灭边界。仍需保证：删除的是真实节点；哨兵不进入 Map；清空后 \code{head.next=tail} 且 \code{tail.prev=head}。

\section{生成性例三：LFU 的二维顺序}

LFU 的淘汰序不是单轴，而是字典序：先比较 frequency，再在同频率内按 recency 淘汰最旧者。
\[
\text{Eviction Priority}=(\text{frequency ascending},\ \text{recency ascending}).
\]
因此需要：
\begin{itemize}
  \item \code{keyMap}: key $\to$ value, frequency 与节点位置；
  \item \code{freqMap}: frequency $\to$ 该频率下按 recency 排列的集合/链表；
  \item \code{minFreq}: 当前非空最小频率。
\end{itemize}
一次命中要把 key 从频率 $f$ 的桶删除、必要时删除空桶、更新 \code{minFreq}，再把它插入 $f+1$ 桶的最近端。插入新 key 后 \code{minFreq=1}。如果同频率容器只是无序集合，就无法兑现“频率相同淘汰最久未使用”的合同。

\begin{warnbox}{“所有操作 $O(1)$”的隐藏前提}
必须选用支持已知节点常数删除、并保持桶内顺序的结构。若每次为了找到同频率最旧项而扫描，淘汰不再是 $O(1)$；若用平衡树维护频率，界会变成 $O(\log n)$。复杂度声明必须与具体表示一致。
\end{warnbox}

\section{动态排名：Top-$K$、Quickselect 与双堆中位数}

\subsection{Top-$K$ 不是“看到 K 就用堆”}
选择取决于数据是否流式、是否允许改写、是否只查询一次：
\begin{center}\small
\begin{tabularx}{\linewidth}{L{3.0cm}L{3.2cm}YY}
\toprule
场景 & 候选 & 成本 & 边界\\
\midrule
流式维护最大 $K$ 个 & 大小为 $K$ 的最小堆 & $O(N\log K)$，空间 $O(K)$ & 可随到随更\\
静态数组单次第 $K$ 大 & Quickselect & 平均 $O(N)$，最坏 $O(N^2)$ & 会改写数组；随机 pivot 降风险\\
需要完整有序输出 & 排序 & $O(N\log N)$ & 多做了全序但满足输出\\
值域小且已知 & 计数/桶 & $O(N+R)$ & 成本绑定值域 $R$\\
\bottomrule
\end{tabularx}
\end{center}
“最大 $K$ 个”用最小堆，是因为堆顶承担当前候选集合的淘汰门槛；新元素不超过门槛即可丢弃，超过门槛则替换。若求最小 $K$ 个，方向反转为最大堆。

\subsection{局部 Top-$K$ 可以合成全局 Top-$K$}
若数据被分成互不重叠的若干分片，每片只保留自己的 Top-$K$，全局 Top-$K$ 一定包含在这些局部候选的并集中：任何未进入某片 Top-$K$ 的元素，在同片至少有 $K$ 个元素不小于它，因此不可能进入全局 Top-$K$。这给出分布式/分片合并的正确性，而不是只凭“MapReduce 常这样做”。

\subsection{频率 Top-$K$：先聚合，再维护淘汰门槛}
“出现次数最多的 $K$ 个”不是直接对原始记录做大小堆：第一阶段必须用 Hash/分片 Map 将相同 key 聚合成 \code{(key,count)}，第二阶段才在频次流上维护大小为 $K$ 的最小堆。堆顶是当前候选的最低频次；新计数超过门槛才替换。若频次相同，必须在堆比较器中固定 key 的字典序或首次出现时间，否则“第 $K$ 个”没有确定合同。全量内存方案为 $O(N)$ 计数加 $O(U\log K)$ 堆维护（$U$ 为不同 key 数）；分片方案先在每片聚合并保留局部 Top-$K$，再对局部候选重新聚合/合并，正确性依赖同 key 必须进入同片或在第二阶段再次合并。

Count-Min Sketch 可以用较小空间给出频次上界并筛出 heavy-hitter 候选，但哈希碰撞会造成高估；它不能单独承担精确 Top-$K$ 的最终裁决。可行的工程合同是“Sketch 预筛 + 精确 Map 二次确认”，并把误差界、概率和 ties 写进 API。

\subsection{动态中位数：分区不变量比两个堆的名字重要}
用最大堆保存较小一半、最小堆保存较大一半，维护：
\[
\boxed{
\max(L)\le \min(R),\qquad |L|-|R|\in\{0,1\}.}
\]
插入时可先进入左堆，再把左堆最大值转移到右堆，若右堆更大再转回一个；这种固定修复顺序同时恢复分区有序和规模平衡。奇数个元素时中位数为 $\max(L)$，偶数时为两堆顶平均值；若数值可能溢出，求和前必须提升类型。

\subsection{关注流信息流：Hash 定位 + 时间线 + $K$ 路归并}
“返回自己和所有关注者的最近十条动态”是一个典型组合接口。Hash 保存 \code{user -> follows} 和 \code{author -> timeline}，每个作者的时间线保持按时间倒序；查询时把自己与关注者的时间线头元素放入最大堆，每次弹出最新一条，再只推进同一来源的下一条。若来源数为 $F$、输出条数为 $K$，查询成本约为 $O(F+K\log F)$，而不是把所有历史动态收集后全量排序；空间为堆中的 $O(F)$ 个 frontier。实现必须明确：用户自身是否总被包含、同一时间戳的 tie-break、follow/unfollow 后对后续查询的生效时刻、已删除动态是否跳过，以及分页游标如何避免重复返回。

全量收集再排序仍可作为小规模、低频查询的正确基线；当历史长度远大于 $K$ 时，$K$ 路归并才兑现了“成本绑定来源数与输出量”的承诺。这里的图关系、Hash 定位、时间序列和堆 frontier 各自有 Owner，I01 只拥有它们如何组合及跨视图一致性。

\section{规模升级：从内存结构到外存与分布式处理}

\subsection{第一问不是“用什么算法”，而是“状态放得下吗”}
对 $N$ 条记录，先估算：
\[
\text{Working Set}
\approx N\times(\text{payload}+\text{index overhead}+\text{allocator/runtime overhead}).
\]
不能只用源文件字节数判断 HashSet/HashMap 是否能装入内存；指针、桶、对象头、装填率和字符串对象可能把内存放大数倍。安全分片数应留出输入缓冲、输出缓冲、运行时与倾斜余量：
\[
K\ge
\left\lceil
\frac{\text{Estimated Working Set}}
     {\text{Usable Memory per Worker}}
\right\rceil.
\]

\subsection{四级规模模型}
\begin{enumerate}
  \item \kw{Fits in memory：}直接使用 Hash、Heap、排序或有序索引；
  \item \kw{Streaming：}输入顺序扫描，只保存固定或次线性摘要，如 Top-$K$ 小堆；
  \item \kw{Partitioned external：}按 key/range 分片，使每片可装入内存后逐片处理；
  \item \kw{Distributed：}各节点处理局部状态，随后按可证明的合并接口归约。
\end{enumerate}
规模升级不是自动优化。每升一级都会增加临时文件、序列化、网络、失败恢复与倾斜处理成本。

\subsection{Hash 分片的正确性条件}
去重、频率统计与交集问题常依赖：
\[
x=y\implies h(x)\bmod K=h(y)\bmod K.
\]
因此相同 key 必须进入同一编号分片。两个文件求交集时必须使用相同规范化规则、Hash 函数、种子与分片数，再逐对处理对应分片。Hash 冲突只让不同 key 落入同片，不会破坏完整性；最终相等判断仍必须比较真实 key，不能把 Hash 值相等当成对象相等。

数据倾斜可能让某片仍超内存，正确恢复是对该片使用新种子/更多位再次分片，或切换范围分区；“平均每片大小足够”不是最坏保证。

\subsection{存在性、去重、交集与频率是不同合同}
\begin{center}\small
\begin{tabularx}{\linewidth}{L{2.5cm}L{3.1cm}YY}
\toprule
任务 & 精确方案主干 & 可选压缩 & 关键边界\\
\midrule
存在性查询 & HashSet / bitmap & Bloom filter & Bloom 有假阳性，不能单独证明存在\\
去重 & Hash 分片后逐片 Set & 排序后相邻去重 & 输出顺序、重复定义、规范化\\
两文件交集 & 同分片后 Set probe & 外部排序 + 双指针 & 两侧必须同分片规则\\
频率统计 & 分片后 Map 计数 & Count-Min Sketch & 近似误差与二次确认\\
Top-$K$ & 固定小堆 & 近似 heavy hitter & ties、顺序与精确性\\
\bottomrule
\end{tabularx}
\end{center}
位图要求 key 能映射到可接受大小的稠密整数域；Bloom filter 用多个位位置换空间，允许“可能存在/一定不存在”，不允许把“可能存在”直接当最终精确答案。

\section{外部排序与范围分桶：瓶颈必须测量}

\subsection{外部归并的 Owner 组合}
当数据不能装入内存且要求全序，主线为：
\[
\text{Run Generation}
\to \text{Internal Sort per Run}
\to \text{Buffered }K\text{-way Merge}
\to \text{Ordered Output}.
\]
多路归并用最小堆保存各 run 当前头元素，总代价约为 $O(N\log K)$ 次候选比较，并伴随顺序块 I/O。内存预算必须同时容纳 $K$ 个输入缓冲、输出缓冲、堆节点和解析状态；不能把全部内存都分给 run 数据。

\subsection{桶排序的成立条件}
范围分桶只有在 key 域可分、桶序能推出全局序、且每桶可被继续处理时才成立。数据分布不均会制造巨大热桶；此时需要采样分位点、递归细分或回退外部归并。按桶编号拼接只在“任意前桶 key 不大于任意后桶 key”时成立，Hash 分片没有这个顺序性质，不能直接拼接成排序结果。

\subsection{CPU、I/O、内存和串行尾部没有固定赢家}
同样是外部处理，不同 workload 的瓶颈可能相反：
\begin{itemize}
  \item 每条记录只解析整数并维护 $K=100$ 的堆时，磁盘吞吐可能占主导；
  \item 每条记录要解析 JSON、提取字段并做昂贵比较时，CPU 可能占主导；
  \item 多线程同时加载大批对象时，内存与 GC 压力可能占主导；
  \item 局部排序已经并行，最终单路/少路归并仍可能形成串行尾部。
\end{itemize}
因此“外部排序一定 I/O bound”与“多线程一定更快”都是错误的无条件规则。先测每阶段的读取、解析、比较、写出、等待、峰值内存与并行度，再决定扩大 batch、增加线程、减少对象化还是改变算法。

\begin{warnbox}{实验数据不是通用定律}
某次 10GB JSON 实验中的 CPU 利用率、耗时、JVM 参数和 GC 行为只能说明该硬件、格式、实现与参数下的瓶颈。Canonical 模型吸收“分阶段测量、定位主导成本和串行比例”的方法，不把具体秒数或百分比晋升为普遍事实。
\end{warnbox}

\section{完整决策表与贯穿案例}

\subsection{七步决策表}
\begin{center}\small
\begin{tabularx}{\linewidth}{L{2.5cm}YY}
\toprule
步骤 & 必须写出 & 验证问题\\
\midrule
Workload & 操作种类、频率、规模、分布、资源 & 是否遗漏主操作或阶段？\\
API Contract & 输入、输出、修改、失败语义、成本界 & “平均/最坏/均摊/近似”是否说清？\\
Capabilities & 定位、顺序、极值、随机、前缀、外存 & 每项能力真的由某 Owner 兑现吗？\\
Candidates & 至少一个正确基线和一个改进候选 & 每个候选满足语义合同吗？\\
Invariants & 单结构与跨结构同步条件 & 每次更新后怎样恢复？\\
Cost Vector & 时间、空间、比较、移动、I/O、通信 & 成本绑定哪个输入和资源层？\\
Choice / Boundary & 取舍、假设、失败与升级路径 & 改变 workload 后选择会否反转？\\
\bottomrule
\end{tabularx}
\end{center}

\subsection{贯穿例子一：日志索引}
任务 A：百万条记录，绝大多数操作是精确 \code{user\_id} 查找，偶尔插入，空间可扩展。候选为 Hash 与平衡树；Hash 以装填和冲突管理换取平均精确查询，平衡树提供最坏对数级和有序遍历。任务 B 增加时间范围扫描后，Hash 需额外维护排序结构，平衡树或 B+ 结构的选择理由增强。结构机制分别由 DS09/DS10 Own，比较接口由 DS-B02 Own。

\subsection{贯穿例子二：最近十条动态}
若每个用户维护按时间递增的动态序列，查询某用户及其关注者的最近十条，直接收集全部动态再排序是正确基线，但成本随历史总量增长。改进方案让每个来源只暴露自己的尾部迭代器，并用大小等于来源数的最大堆做 $K$ 路合并；每弹出一条，只把同一来源的前一条加入候选。于是成本绑定关注来源数与输出条数，而不是全部历史。

这里 Hash 只负责 \code{user -> follows/timeline} 定位，时间序列负责来源内有序，Heap 负责跨来源 frontier。若关注关系、删除语义或分页游标改变，必须重新审查状态与一致性；不能把演示代码中“收集后全排序”的基线冒充最终伸缩方案。

\section{反证、失败边界与出口}

\begin{center}\small
\begin{tabularx}{\linewidth}{L{3.1cm}YY}
\toprule
危险结论 & 最小反例/追问 & 正确出口\\
\midrule
Hash 总是最快 & 范围查询、碰撞最坏界、重散列、顺序输出 & 回 DS09/DS-B02 比较\\
平衡树总是最好 & 只需精确查找且更新频繁，是否值得维护全序 & 允许条件化选择\\
看到 Top-$K$ 就用堆 & 静态数组单次查询且允许改写 & 比较 Quickselect/排序\\
Hash 分片后可直接排序拼接 & Hash 分片无范围顺序 & 使用范围桶或外部归并\\
多线程一定加速 & 文件读取串行、锁竞争、GC、最终归并 & 先做阶段测量\\
Bloom 说存在就一定存在 & 假阳性 & 二次精确确认\\
组合结构各自正确即可 & 两个索引可互相失配 & 写跨结构同步不变量\\
\bottomrule
\end{tabularx}
\end{center}

\section{验证协议：不仅验证结果，也验证状态}
组合结构至少使用以下测试：
\begin{enumerate}
  \item 空结构、容量为零、单元素、重复插入、删除不存在 key；
  \item 删除数组末元素与非末元素，随后再次定位被搬移元素；
  \item LRU 更新旧 key、连续命中、容量淘汰与重新插入；
  \item LFU 同频率 tie-break、桶清空、\code{minFreq} 跳转与容量为零；
  \item Top-$K$ 的 $K=1$、$K=N$、重复值与 ties；
  \item 双堆中位数的递增、递减、重复、奇偶长度与数值溢出；
  \item 分片后的全量守恒：输入记录数等于各分片记录数之和；
  \item 相同 key 必须同片，逐片结果合并后与小规模内存基准一致；
  \item 人工制造倾斜片、短读写、空 run 与最后不足一个 batch 的尾块。
\end{enumerate}
可在小规模输入上同时运行“慢但显然正确”的基线与优化方案，比较输出集合/序列；这是验证组合正确性的独立信息，不应只用优化代码自证。

\section{做题控制协议与压缩复原}
草稿固定写九行：
\begin{enumerate}
  \item API 与失败语义；
  \item 操作频率与规模；
  \item 每项操作要求的能力和复杂度口径；
  \item 正确基线；
  \item 基线重复工作/资源瓶颈；
  \item 候选结构与能力分工；
  \item 跨结构同步不变量；
  \item 成本向量、假设与边界；
  \item 最小反例和规模升级路线。
\end{enumerate}
最后删除没有改变选择的装饰细节；保留能让下一道陌生题复用的接口句。

\begin{corebox}{压缩复原：九问}
真正昂贵的操作是什么？成本要求是平均、最坏还是均摊？哪个结构提前保存了所需信息？单一结构缺哪项能力？多个结构怎样指向同一份逻辑事实？一次修改要同步哪些视图？数据是否装得下？若装不下，按什么不变量分片并怎样合并？改变 workload 后选择是否反转？
\end{corebox}

\end{document}
